\subsubsection{Heavy-Light Decomposition (HLD)}
\begin{lstlisting}[style=mystyle, language=C++]
// TC: O(N) build, O(log^2 N) por query en camino entre u y v
// usa: descomponer arbol en cadenas pesadas para queries de camino con SegTree
#include <bits/stdc++.h>
using namespace std;

vector<int> heavy, sz, parent, depth, pos, head;
int t = -1;

void dfs1(int node, int par, int d, vector<int> adj[]){
  parent[node] = par;
  depth[node] = d;
  sz[node] = 1;
  int mx = 0;
  for(auto adjN : adj[node]) if(adjN != par){
    dfs1(adjN, node, d + 1, adj);
    sz[node] += sz[adjN];
    if(sz[adjN] > mx){
      mx = sz[adjN];
      heavy[node] = adjN;
    }
  }
}

void dfs2(int node, int h, vector<int> adj[]){
  head[node] = h;
  pos[node] = ++t;
  if(heavy[node] != -1) dfs2(heavy[node], h, adj);
  for(auto adjN : adj[node]){
    if(adjN == parent[node] || adjN == heavy[node]) continue;
    dfs2(adjN, adjN, adj);
  }
}

void init_hld(int root, int n, vector<int> adj[]){
  parent.assign(n, -1);
  heavy.assign(n, -1);
  sz.assign(n, 1);
  depth.assign(n, 0);
  pos.assign(n, 0);
  head.assign(n, 0);
  t = -1;
  dfs1(root, -1, 0, adj);
  dfs2(root, root, adj);
}

// Procesa los rangos [pos[head[u]], pos[u]] en el SegTree
// Ejemplo para queries en camino u <-> v (suma, max, etc.):
int query_path(int u, int v){
  int res = 0; // elemento neutro de la operacion
  while(head[u] != head[v]){
    if(depth[head[u]] < depth[head[v]]) swap(u, v);
    // res = combine(res, query_segtree(pos[head[u]], pos[u]));
    u = parent[head[u]];
  }
  if(depth[u] > depth[v]) swap(u, v);
  // res = combine(res, query_segtree(pos[u], pos[v]));
  return res;
}
\end{lstlisting}
